\documentclass[12pt]{article}

\usepackage[margin=1in]{geometry}
\usepackage{setspace}
\doublespacing
\usepackage{amsmath,amssymb}
\usepackage{booktabs}
\usepackage{threeparttable}
\usepackage{caption}
\usepackage{array}
\usepackage{longtable}
\usepackage{natbib}
\bibpunct{(}{)}{;}{a}{,}{,}
\usepackage{hyperref}
\hypersetup{colorlinks=true, linkcolor=black, citecolor=black, urlcolor=black}
\usepackage{microtype}
\renewcommand{\thetable}{A\arabic{table}}

\title{\textbf{Online Appendix}\\
\large The Market Pricing of Proprietary AI Model Capability}
\author{\begin{tabular}{c}
Zhuo Chen\textsuperscript{a}\quad Ruishuang Lu\textsuperscript{a}\\[0.5em]
{\footnotesize \textsuperscript{a} Center of Economic Research, Shandong University}
\end{tabular}}
\date{\small This version: June 2026}

\begin{document}
\maketitle

\section{Sample Construction}
\label{app:sample}

The starting point is a set of 60 major AI model release events between January
2024 and March 2026. The event screen requires an official model or model-family
release, a confirmed public release date, public access or benchmark visibility,
and an economic link to at least one AI-related listed firm. The clean
firm-event panel has 5,160 observations and 86 firms. The main capability sample
has 47 releases with nonmissing Artificial Analysis (AA) Intelligence Index
scores and 3,780 complete firm-event observations. The closed-source sample has
36 events and 2,899 observations. The open-weight sample has 11 events and 881
observations.

\begin{table}[htbp]
\centering
\caption{Sample summary}
\label{tab:a_sample}
\begin{threeparttable}
\begin{tabular}{lrr}
\toprule
Sample component & Observations & Events \\
\midrule
Clean firm-event panel & 5,160 & 60 \\
Main capability sample & 3,780 & 47 \\
Closed-source capability sample & 2,899 & 36 \\
Open-weight capability sample & 881 & 11 \\
Pre-event CAR sample & 3,762 & 47 \\
FF3 CAR sample & 3,782 & 47 \\
US-traded ticker screen & 2,857 & 47 \\
\bottomrule
\end{tabular}
\begin{tablenotes}[flushleft]\footnotesize
\item The main capability sample requires nonmissing AA Intelligence Index,
CAR$[0,{+}20]$, firm controls, year, event identifier, and firm identifier.
The US-traded ticker screen uses a rule-based ticker filter and should be read as
a robustness check rather than a formal exchange-country classification.
\end{tablenotes}
\end{threeparttable}
\end{table}

\section{Variable Definitions}
\label{app:variables}

The dependent variable is a cumulative abnormal return. The market-model CAR uses
a firm-level market model fitted before the release date. The FF3 CAR uses the
Fama-French three-factor model. The main window is $[0,{+}20]$ trading days. The
pre-event CAR uses the pre-release window available in the cleaned panel. The main
capability measure is the AA Intelligence Index. The open-weight indicator equals
one when the model is released under an open-weight or open-source regime.

Firm-event relationship variables are binary indicators. They include owner,
investor, cloud, business upstream, real upstream, business downstream, real
downstream, and competitor. The relationship variables were merged into the
cleaned panel using event and firm identifiers. They are used for heterogeneity
and appendix checks, not for the baseline result.

\section{Main Results with Small-Cluster Inference}
\label{app:smallcluster}

Table~\ref{tab:a_core} reports the core estimates with CR0, CR2, and wild cluster
bootstrap inference. The closed-source CAR$[0,{+}20]$ result is the most stable
estimate. The full-sample slope and open-weight interaction are weaker under the
most conservative inference and should be read as supporting evidence.

\begin{table}[htbp]
\centering
\caption{Core estimates with event-cluster inference}
\label{tab:a_core}
\begin{threeparttable}
\begin{tabular}{lrrrrrr}
\toprule
Specification & N & Events & Coef. & SE & CR2 $p$ & Wild $p$ \\
\midrule
Full sample, CAR$[0,{+}20]$ & 3,780 & 47 & 0.00152 & 0.00067 & 0.054 & 0.053 \\
Closed-source, CAR$[0,{+}20]$ & 2,899 & 36 & 0.00232 & 0.00062 & 0.007 & 0.008 \\
Interaction, closed-source slope & 3,780 & 47 & 0.00226 & 0.00060 & 0.005 & 0.008 \\
Interaction, open-weight term & 3,780 & 47 & -0.00373 & 0.00120 & 0.048 & 0.086 \\
Investor interaction & 3,780 & 47 & -0.00272 & 0.00098 & 0.050 & 0.126 \\
Cloud interaction & 3,780 & 47 & -0.00213 & 0.00104 & 0.100 & 0.137 \\
Owner interaction, CAR$[0,{+}1]$ & 3,763 & 47 & -0.00178 & 0.00072 & 0.059 & 0.034 \\
Owner interaction, CAR$[0,{+}2]$ & 3,768 & 47 & -0.00177 & 0.00070 & 0.052 & 0.046 \\
\bottomrule
\end{tabular}
\begin{tablenotes}[flushleft]\footnotesize
\item The coefficient is on the AA Intelligence Index or the stated interaction.
Wild $p$-values use 4,999 Rademacher draws clustered by event. The investor,
cloud, and owner rows are exploratory relationship checks.
\end{tablenotes}
\end{threeparttable}
\end{table}

\section{Firm Fixed Effects, Alternative Returns, and Pre-Event CARs}
\label{app:robust}

Table~\ref{tab:a_robust} reports robustness checks that address firm composition,
alternative abnormal-return models, pre-event pricing, and the exchange-country
mix in the firm panel. The closed-source coefficient remains positive in all
post-event specifications. The pre-event coefficient is not statistically
significant.

\begin{table}[htbp]
\centering
\caption{Robustness checks for the capability slope}
\label{tab:a_robust}
\begin{threeparttable}
\small
\begin{tabular}{lrrrrr}
\toprule
Specification & N & Events & Coef. & SE & $p$ \\
\midrule
All, firm fixed effects, event cluster & 3,780 & 47 & 0.00154 & 0.00067 & 0.026 \\
All, firm fixed effects, two-way cluster & 3,780 & 47 & 0.00154 & 0.00077 & 0.052 \\
Closed-source, firm fixed effects, event cluster & 2,899 & 36 & 0.00226 & 0.00063 & 0.001 \\
Closed-source, firm fixed effects, two-way cluster & 2,899 & 36 & 0.00226 & 0.00087 & 0.014 \\
All, FF3 CAR$[0,{+}20]$ & 3,782 & 47 & 0.00102 & 0.00044 & 0.025 \\
Closed-source, FF3 CAR$[0,{+}20]$ & 2,901 & 36 & 0.00118 & 0.00051 & 0.027 \\
All, pre-event CAR & 3,762 & 47 & 0.00016 & 0.00028 & 0.576 \\
Closed-source, pre-event CAR & 2,884 & 36 & 0.00047 & 0.00030 & 0.119 \\
All, US-traded ticker screen & 2,857 & 47 & 0.00193 & 0.00080 & 0.019 \\
Closed-source, US-traded ticker screen & 2,191 & 36 & 0.00271 & 0.00072 & 0.001 \\
\bottomrule
\end{tabular}
\begin{tablenotes}[flushleft]\footnotesize
\item The table reports the coefficient on the AA Intelligence Index. Two-way
clustering clusters by event and firm. The US-traded ticker screen is based on
ticker format and is used only as a robustness check.
\end{tablenotes}
\end{threeparttable}
\end{table}

\section{Event-Level Aggregation}
\label{app:eventlevel}

The main panel repeats the event-level capability score across firms. Table
\ref{tab:a_eventlevel} collapses the data to the event level by averaging CARs
within each release. The closed-source result survives when all firms are used in
the equal-weighted event mean. The result is weaker when only firms with explicit
relationship indicators are averaged.

\begin{table}[htbp]
\centering
\caption{Event-level aggregation}
\label{tab:a_eventlevel}
\begin{threeparttable}
\begin{tabular}{lrrrr}
\toprule
Event-level outcome & Events & Coef. & SE & $p$ \\
\midrule
All firms, all releases & 47 & 0.00099 & 0.00062 & 0.120 \\
All firms, closed-source releases & 36 & 0.00172 & 0.00067 & 0.016 \\
Relationship firms, all releases & 47 & 0.00061 & 0.00078 & 0.437 \\
Relationship firms, closed-source releases & 36 & 0.00142 & 0.00103 & 0.178 \\
\bottomrule
\end{tabular}
\begin{tablenotes}[flushleft]\footnotesize
\item Each row estimates an event-level regression in which the dependent
variable is an equal-weighted mean CAR$[0,{+}20]$ within the release event.
Relationship firms are firm-event pairs with at least one explicit relationship
indicator.
\end{tablenotes}
\end{threeparttable}
\end{table}

\section{Leave-One-Creator-Out}
\label{app:loo}

Table~\ref{tab:a_loo} excludes major model creators one at a time. The
closed-source coefficient remains positive and statistically significant in each
reported exclusion.

\begin{table}[htbp]
\centering
\caption{Leave-one-creator-out, closed-source releases}
\label{tab:a_loo}
\begin{threeparttable}
\begin{tabular}{lrrrrr}
\toprule
Excluded creator & N & Events & Coef. & SE & $p$ \\
\midrule
Google & 1,927 & 24 & 0.00204 & 0.00077 & 0.015 \\
OpenAI & 2,010 & 25 & 0.00426 & 0.00069 & 0.000 \\
Anthropic & 2,272 & 28 & 0.00198 & 0.00062 & 0.004 \\
Alibaba & 2,820 & 35 & 0.00215 & 0.00061 & 0.001 \\
Meta & 2,816 & 35 & 0.00235 & 0.00061 & 0.001 \\
xAI & 2,733 & 34 & 0.00225 & 0.00066 & 0.002 \\
Microsoft & 2,816 & 35 & 0.00207 & 0.00070 & 0.006 \\
\bottomrule
\end{tabular}
\begin{tablenotes}[flushleft]\footnotesize
\item The table reports the closed-source coefficient on the AA Intelligence
Index after excluding all events from the listed creator.
\end{tablenotes}
\end{threeparttable}
\end{table}

\section{Specification-Curve Evidence}
\label{app:specr}

Table~\ref{tab:a_specr_metrics} compares alternative capability measures across
reasonable specifications. The AA Intelligence Index has the strongest overall
pattern among the LLM capability measures. The AA Media Elo measure behaves
differently and is concentrated in short-window responses.

\begin{table}[htbp]
\centering
\caption{Specification-curve summary by capability measure}
\label{tab:a_specr_metrics}
\begin{threeparttable}
\begin{tabular}{lrrrr}
\toprule
Capability measure & Specifications & 5\% significant & Positive & Median coefficient \\
\midrule
AA Intelligence Index & 252 & 23.8\% & 66.3\% & 0.00020 \\
AA Media Elo & 231 & 20.3\% & 62.3\% & 0.00002 \\
AA Coding Index & 252 & 11.5\% & 51.6\% & 0.00001 \\
AA Math Index & 231 & 5.6\% & 26.8\% & -0.00009 \\
\bottomrule
\end{tabular}
\begin{tablenotes}[flushleft]\footnotesize
\item The specification curve varies the event window, control set, subsample,
and year fixed effects. The table reports the share of estimates significant at
5 percent using the corresponding clustered standard error in each specification.
\end{tablenotes}
\end{threeparttable}
\end{table}

\begin{table}[htbp]
\centering
\caption{AA Intelligence Index across event windows}
\label{tab:a_windows}
\begin{tabular}{lrr}
\toprule
Window & 5\% significant & Median coefficient \\
\midrule
CAR$[0,{+}1]$ & 30.6\% & 0.000124 \\
CAR$[0,{+}2]$ & 25.0\% & 0.000091 \\
CAR$[0,{+}3]$ & 22.2\% & 0.000080 \\
CAR$[0,{+}5]$ & 2.8\% & -0.000014 \\
CAR$[0,{+}10]$ & 13.9\% & -0.000247 \\
CAR$[0,{+}15]$ & 25.0\% & 0.000507 \\
CAR$[0,{+}20]$ & 47.2\% & 0.000848 \\
\bottomrule
\end{tabular}
\end{table}

\section{Additional Heterogeneity}
\label{app:heterogeneity}

\subsection{Open-Weight and Closed-Source Specification Patterns}

The long-window open-weight estimates are negative in point estimates, while the
closed-source specifications are positive. Because the open-weight group has only
11 events, the main text treats the interaction result as suggestive mechanism
evidence.

\begin{table}[htbp]
\centering
\caption{Open-weight and closed-source patterns}
\label{tab:a_open_closed}
\begin{tabular}{lrrr}
\toprule
Subsample & 5\% significant & Positive & Median coefficient \\
\midrule
Open-weight releases & 38.1\% & 7.1\% & -0.000599 \\
Closed-source releases & 26.2\% & 88.1\% & 0.000512 \\
\bottomrule
\end{tabular}
\end{table}

\subsection{Owner Firms}

Owner firms display a short-window negative response that dissipates by 10 to 20
trading days. This is an exploratory result because the owner sample is small.

\begin{table}[htbp]
\centering
\caption{Owner-firm short-window response}
\label{tab:a_owner}
\begin{tabular}{lrrr}
\toprule
Window & Coef. & SE & $p$ \\
\midrule
CAR$[0,{+}1]$ & -0.00232 & 0.00068 & 0.003 \\
CAR$[0,{+}2]$ & -0.00287 & 0.00063 & 0.000 \\
CAR$[0,{+}3]$ & -0.00270 & 0.00077 & 0.002 \\
CAR$[0,{+}5]$ & -0.00205 & 0.00102 & 0.057 \\
CAR$[0,{+}10]$ & -0.00080 & 0.00201 & 0.695 \\
CAR$[0,{+}20]$ & 0.00047 & 0.00285 & 0.871 \\
\bottomrule
\end{tabular}
\end{table}

\subsection{Event Tier}

Tier 2 events show a larger capability slope than Tier 1 events. The interaction
between intelligence and Tier 1 is not statistically significant, so we treat
this as a descriptive split.

\begin{table}[htbp]
\centering
\caption{Capability slope by event tier}
\label{tab:a_tier}
\begin{tabular}{lrrrr}
\toprule
Tier & N & Events & CAR$[0,{+}20]$ coefficient & $p$ \\
\midrule
Tier 1 & 1,438 & 18 & 0.00153 & 0.045 \\
Tier 2 & 2,159 & 27 & 0.00374 & 0.000 \\
\bottomrule
\end{tabular}
\end{table}

\subsection{Calendar Year}

The capability slope is not significant in 2024 and turns positive in 2025.
The 2026 estimate is large but imprecise because the sample is small.

\begin{table}[htbp]
\centering
\caption{Capability slope by release year}
\label{tab:a_year}
\begin{tabular}{lrrrr}
\toprule
Year & N & Coef. & SE & $p$ \\
\midrule
2024 & 1,291 & -0.00044 & 0.00060 & 0.474 \\
2025 & 2,183 & 0.00206 & 0.00087 & 0.026 \\
2026 & 306 & 0.01860 & 0.01157 & 0.206 \\
\bottomrule
\end{tabular}
\end{table}

\subsection{Industry}

The strongest industry-level result appears in internet services and
infrastructure. Software and semiconductors show positive but more marginal
patterns.

\begin{table}[htbp]
\centering
\caption{Industry-level capability slopes}
\label{tab:a_industry}
\begin{tabular}{lrrrr}
\toprule
Industry & N & Coef. & CR0 $p$ & CR2 $p$ \\
\midrule
Internet services and infrastructure & 235 & 0.00190 & 0.005 & 0.015 \\
Software & 1,330 & 0.00295 & 0.045 & 0.078 \\
Semiconductors and semiconductor equipment & 606 & 0.00212 & 0.032 & 0.060 \\
Technology hardware, storage, and peripherals & 348 & -0.00083 & 0.177 & 0.210 \\
IT services & 325 & 0.00047 & 0.533 & 0.561 \\
Internet retail & 184 & -0.00004 & 0.947 & 0.950 \\
Other or non-IT & 752 & 0.00027 & 0.567 & 0.600 \\
\bottomrule
\end{tabular}
\end{table}

\subsection{Mag7 and Non-Mag7 Firms}

The Mag7 and non-Mag7 slopes are similar. The interaction between intelligence
and the Mag7 indicator is not statistically significant.

\begin{table}[htbp]
\centering
\caption{Mag7 and non-Mag7 slopes}
\label{tab:a_mag7}
\begin{tabular}{lrrrr}
\toprule
Group & N & Coef. & CR0 $p$ & CR2 $p$ \\
\midrule
Mag7 & 328 & 0.00170 & 0.045 & 0.081 \\
Non-Mag7 & 3,452 & 0.00147 & 0.037 & 0.068 \\
\bottomrule
\end{tabular}
\end{table}

\section{Media Sentiment}
\label{app:sentiment}

FinBERT sentiment is negatively related to CAR$[0,{+}20]$. In joint models, the
sentiment and capability coefficients retain their signs but are marginal under
CR2. The sentiment-capability interaction is not statistically significant.

\begin{table}[htbp]
\centering
\caption{Media sentiment and CAR$[0,{+}20]$}
\label{tab:a_sentiment}
\begin{tabular}{lrrr}
\toprule
Specification & Coef. & CR0 $p$ & CR2 $p$ \\
\midrule
FinBERT sentiment, $\pm 5$ days & -0.081 & 0.000 & 0.003 \\
FinBERT sentiment, $\pm 20$ days & -0.096 & 0.003 & 0.015 \\
Sentiment in joint model & -0.061 & 0.024 & 0.057 \\
Intelligence in joint model & 0.00158 & 0.035 & 0.087 \\
\bottomrule
\end{tabular}
\end{table}

\section{Event-Window Overlap}
\label{app:overlap}

AI model releases are clustered in calendar time. Among the 47 intelligence
events, only four have no other intelligence event within 20 calendar days.
The median event has three other events within that window. Because a strict
non-overlap sample would be too small, the main paper keeps the full capability
sample and reports this diagnostic.

\begin{table}[htbp]
\centering
\caption{Overlap within 20 calendar days}
\label{tab:a_overlap}
\begin{tabular}{lr}
\toprule
Statistic & Other intelligence events within 20 days \\
\midrule
Minimum & 0 \\
25th percentile & 2 \\
Median & 3 \\
Mean & 2.68 \\
75th percentile & 4 \\
Maximum & 6 \\
Events with no overlap & 4 of 47 \\
\bottomrule
\end{tabular}
\end{table}

\section{Null and Exploratory Results}
\label{app:null}

Several results are informative but should not be interpreted as main findings.
The competitor subsample has a small and statistically insignificant capability
slope. Real downstream firms have a larger point estimate but too few
observations for a precise estimate. Investor and cloud interactions are negative
but do not survive the most conservative inference.

\begin{table}[htbp]
\centering
\caption{Exploratory and null relationship results}
\label{tab:a_exploratory}
\begin{tabular}{lrrrr}
\toprule
Result & N & Events & Coef. & $p$ \\
\midrule
Competitor subsample & 366 & 47 & 0.00052 & 0.352 \\
Real downstream subsample & 62 & 20 & 0.00200 & 0.235 \\
Investor interaction, wild bootstrap & 3,780 & 47 & -0.00272 & 0.126 \\
Cloud interaction, wild bootstrap & 3,780 & 47 & -0.00213 & 0.137 \\
\bottomrule
\end{tabular}
\end{table}

\end{document}
